% Generated by Sphinx.
\def\sphinxdocclass{report}
\documentclass[letterpaper,10pt,english]{sphinxmanual}
\usepackage[utf8]{inputenc}
\DeclareUnicodeCharacter{00A0}{\nobreakspace}
\usepackage{cmap}
\usepackage[T1]{fontenc}
\usepackage{babel}
\usepackage{times}
\usepackage[Bjarne]{fncychap}
\usepackage{longtable}
\usepackage{sphinx}
\usepackage{multirow}

\usepackage{dsfont}
\usepackage{slashed}
\usepackage{yfonts}
\usepackage{mathrsfs}
\def\degrees{^\circ}
\def\d{{\rm d}}

\def\sign{\mathop{\mathrm{sign}}}
\def\L{{\mathcal L}}
\def\H{{\mathcal H}}
\def\M{{\mathcal M}}
\def\matrix{}
\def\fslash#1{#1 \!\!\!/}
\def\F{{\bf F}}
\def\R{{\bf R}}
\def\J{{\bf J}}
\def\x{{\bf x}}
\def\y{{\bf y}}
\def\h{{\rm h}}
\def\a{{\rm a}}
\newcommand{\bfx}{\mbox{\boldmath $x$}}
\newcommand{\bfy}{\mbox{\boldmath $y$}}
\newcommand{\bfz}{\mbox{\boldmath $z$}}
\newcommand{\bfv}{\mbox{\boldmath $v$}}
\newcommand{\bfu}{\mbox{\boldmath $u$}}
\newcommand{\bfF}{\mbox{\boldmath $F$}}
\newcommand{\bfJ}{\mbox{\boldmath $J$}}
\newcommand{\bfU}{\mbox{\boldmath $U$}}
\newcommand{\bfY}{\mbox{\boldmath $Y$}}
\newcommand{\bfR}{\mbox{\boldmath $R$}}
\newcommand{\bfg}{\mbox{\boldmath $g$}}
\newcommand{\bfc}{\mbox{\boldmath $c$}}
\newcommand{\bfxi}{\mbox{\boldmath $\xi$}}

\newcommand{\bra}[1]{\left\langle #1\right|}
\newcommand{\ket}[1]{\left| #1\right\rangle}
\newcommand{\braket}[2]{\langle #1 \mid #2 \rangle}
\newcommand{\avg}[1]{\left< #1 \right>}


%\def\back{\!\!\!\!\!\!\!\!\!\!}
\def\back{}
\def\col#1#2{\left(\matrix{#1#2}\right)}
\def\row#1#2{\left(\matrix{#1#2}\right)}
\def\mat#1{\begin{pmatrix}#1\end{pmatrix}}
\def\matd#1#2{\left(\matrix{#1\back0\cr0\back#2}\right)}
\def\p#1#2{{\partial#1\over\partial#2}}
\def\cg#1#2#3#4#5#6{({#1},\,{#2},\,{#3},\,{#4}\,|\,{#5},\,{#6})}
\def\half{{\textstyle{1\over2}}}
\def\jsym#1#2#3#4#5#6{\left\{\matrix{
{#1}{#2}{#3}
{#4}{#5}{#6}
}\right\}}
\def\diag{\hbox{diag}}

\font\dsrom=dsrom10
\def\one{\hbox{\dsrom 1}}

\def\res{\mathop{\mathrm{Res}}}

\def\mathnot#1{\text{"$#1$"}}


%See Character Table for cmmib10:
%http://www.math.union.edu/~dpvc/jsmath/download/extra-fonts/cmmib10/cmmib10.html
\font\mib=cmmib10
\def\balpha{\hbox{\mib\char"0B}}
\def\bbeta{\hbox{\mib\char"0C}}
\def\bgamma{\hbox{\mib\char"0D}}
\def\bdelta{\hbox{\mib\char"0E}}
\def\bepsilon{\hbox{\mib\char"0F}}
\def\bzeta{\hbox{\mib\char"10}}
\def\boldeta{\hbox{\mib\char"11}}
\def\btheta{\hbox{\mib\char"12}}
\def\biota{\hbox{\mib\char"13}}
\def\bkappa{\hbox{\mib\char"14}}
\def\blambda{\hbox{\mib\char"15}}
\def\bmu{\hbox{\mib\char"16}}
\def\bnu{\hbox{\mib\char"17}}
\def\bxi{\hbox{\mib\char"18}}
\def\bpi{\hbox{\mib\char"19}}
\def\brho{\hbox{\mib\char"1A}}
\def\bsigma{\hbox{\mib\char"1B}}
\def\btau{\hbox{\mib\char"1C}}
\def\bupsilon{\hbox{\mib\char"1D}}
\def\bphi{\hbox{\mib\char"1E}}
\def\bchi{\hbox{\mib\char"1F}}
\def\bpsi{\hbox{\mib\char"20}}
\def\bomega{\hbox{\mib\char"21}}

\def\bvarepsilon{\hbox{\mib\char"22}}
\def\bvartheta{\hbox{\mib\char"23}}
\def\bvarpi{\hbox{\mib\char"24}}
\def\bvarrho{\hbox{\mib\char"25}}
\def\bvarphi{\hbox{\mib\char"27}}

%how to use:
%$$\alpha\balpha$$
%$$\beta\bbeta$$
%$$\gamma\bgamma$$
%$$\delta\bdelta$$
%$$\epsilon\bepsilon$$
%$$\zeta\bzeta$$
%$$\eta\boldeta$$
%$$\theta\btheta$$
%$$\iota\biota$$
%$$\kappa\bkappa$$
%$$\lambda\blambda$$
%$$\mu\bmu$$
%$$\nu\bnu$$
%$$\xi\bxi$$
%$$\pi\bpi$$
%$$\rho\brho$$
%$$\sigma\bsigma$$
%$$\tau\btau$$
%$$\upsilon\bupsilon$$
%$$\phi\bphi$$
%$$\chi\bchi$$
%$$\psi\bpsi$$
%$$\omega\bomega$$
%
%$$\varepsilon\bvarepsilon$$
%$$\vartheta\bvartheta$$
%$$\varpi\bvarpi$$
%$$\varrho\bvarrho$$
%$$\varphi\bvarphi$$

%small font
\font\mibsmall=cmmib7
\def\bsigmasmall{\hbox{\mibsmall\char"1B}}

\def\Tr{\hbox{Tr}\,}
\def\Arg{\hbox{Arg}}
\def\atan{\hbox{atan}}


\title{Quantum Notes}
\date{February 12, 2014}
\release{1.38}
\author{Lei Ma}
\newcommand{\sphinxlogo}{}
\renewcommand{\releasename}{Release}
\makeindex

\makeatletter
\def\PYG@reset{\let\PYG@it=\relax \let\PYG@bf=\relax%
    \let\PYG@ul=\relax \let\PYG@tc=\relax%
    \let\PYG@bc=\relax \let\PYG@ff=\relax}
\def\PYG@tok#1{\csname PYG@tok@#1\endcsname}
\def\PYG@toks#1+{\ifx\relax#1\empty\else%
    \PYG@tok{#1}\expandafter\PYG@toks\fi}
\def\PYG@do#1{\PYG@bc{\PYG@tc{\PYG@ul{%
    \PYG@it{\PYG@bf{\PYG@ff{#1}}}}}}}
\def\PYG#1#2{\PYG@reset\PYG@toks#1+\relax+\PYG@do{#2}}

\expandafter\def\csname PYG@tok@gd\endcsname{\def\PYG@tc##1{\textcolor[rgb]{0.63,0.00,0.00}{##1}}}
\expandafter\def\csname PYG@tok@gu\endcsname{\let\PYG@bf=\textbf\def\PYG@tc##1{\textcolor[rgb]{0.50,0.00,0.50}{##1}}}
\expandafter\def\csname PYG@tok@gt\endcsname{\def\PYG@tc##1{\textcolor[rgb]{0.00,0.27,0.87}{##1}}}
\expandafter\def\csname PYG@tok@gs\endcsname{\let\PYG@bf=\textbf}
\expandafter\def\csname PYG@tok@gr\endcsname{\def\PYG@tc##1{\textcolor[rgb]{1.00,0.00,0.00}{##1}}}
\expandafter\def\csname PYG@tok@cm\endcsname{\let\PYG@it=\textit\def\PYG@tc##1{\textcolor[rgb]{0.25,0.50,0.56}{##1}}}
\expandafter\def\csname PYG@tok@vg\endcsname{\def\PYG@tc##1{\textcolor[rgb]{0.73,0.38,0.84}{##1}}}
\expandafter\def\csname PYG@tok@m\endcsname{\def\PYG@tc##1{\textcolor[rgb]{0.13,0.50,0.31}{##1}}}
\expandafter\def\csname PYG@tok@mh\endcsname{\def\PYG@tc##1{\textcolor[rgb]{0.13,0.50,0.31}{##1}}}
\expandafter\def\csname PYG@tok@cs\endcsname{\def\PYG@tc##1{\textcolor[rgb]{0.25,0.50,0.56}{##1}}\def\PYG@bc##1{\setlength{\fboxsep}{0pt}\colorbox[rgb]{1.00,0.94,0.94}{\strut ##1}}}
\expandafter\def\csname PYG@tok@ge\endcsname{\let\PYG@it=\textit}
\expandafter\def\csname PYG@tok@vc\endcsname{\def\PYG@tc##1{\textcolor[rgb]{0.73,0.38,0.84}{##1}}}
\expandafter\def\csname PYG@tok@il\endcsname{\def\PYG@tc##1{\textcolor[rgb]{0.13,0.50,0.31}{##1}}}
\expandafter\def\csname PYG@tok@go\endcsname{\def\PYG@tc##1{\textcolor[rgb]{0.20,0.20,0.20}{##1}}}
\expandafter\def\csname PYG@tok@cp\endcsname{\def\PYG@tc##1{\textcolor[rgb]{0.00,0.44,0.13}{##1}}}
\expandafter\def\csname PYG@tok@gi\endcsname{\def\PYG@tc##1{\textcolor[rgb]{0.00,0.63,0.00}{##1}}}
\expandafter\def\csname PYG@tok@gh\endcsname{\let\PYG@bf=\textbf\def\PYG@tc##1{\textcolor[rgb]{0.00,0.00,0.50}{##1}}}
\expandafter\def\csname PYG@tok@ni\endcsname{\let\PYG@bf=\textbf\def\PYG@tc##1{\textcolor[rgb]{0.84,0.33,0.22}{##1}}}
\expandafter\def\csname PYG@tok@nl\endcsname{\let\PYG@bf=\textbf\def\PYG@tc##1{\textcolor[rgb]{0.00,0.13,0.44}{##1}}}
\expandafter\def\csname PYG@tok@nn\endcsname{\let\PYG@bf=\textbf\def\PYG@tc##1{\textcolor[rgb]{0.05,0.52,0.71}{##1}}}
\expandafter\def\csname PYG@tok@no\endcsname{\def\PYG@tc##1{\textcolor[rgb]{0.38,0.68,0.84}{##1}}}
\expandafter\def\csname PYG@tok@na\endcsname{\def\PYG@tc##1{\textcolor[rgb]{0.25,0.44,0.63}{##1}}}
\expandafter\def\csname PYG@tok@nb\endcsname{\def\PYG@tc##1{\textcolor[rgb]{0.00,0.44,0.13}{##1}}}
\expandafter\def\csname PYG@tok@nc\endcsname{\let\PYG@bf=\textbf\def\PYG@tc##1{\textcolor[rgb]{0.05,0.52,0.71}{##1}}}
\expandafter\def\csname PYG@tok@nd\endcsname{\let\PYG@bf=\textbf\def\PYG@tc##1{\textcolor[rgb]{0.33,0.33,0.33}{##1}}}
\expandafter\def\csname PYG@tok@ne\endcsname{\def\PYG@tc##1{\textcolor[rgb]{0.00,0.44,0.13}{##1}}}
\expandafter\def\csname PYG@tok@nf\endcsname{\def\PYG@tc##1{\textcolor[rgb]{0.02,0.16,0.49}{##1}}}
\expandafter\def\csname PYG@tok@si\endcsname{\let\PYG@it=\textit\def\PYG@tc##1{\textcolor[rgb]{0.44,0.63,0.82}{##1}}}
\expandafter\def\csname PYG@tok@s2\endcsname{\def\PYG@tc##1{\textcolor[rgb]{0.25,0.44,0.63}{##1}}}
\expandafter\def\csname PYG@tok@vi\endcsname{\def\PYG@tc##1{\textcolor[rgb]{0.73,0.38,0.84}{##1}}}
\expandafter\def\csname PYG@tok@nt\endcsname{\let\PYG@bf=\textbf\def\PYG@tc##1{\textcolor[rgb]{0.02,0.16,0.45}{##1}}}
\expandafter\def\csname PYG@tok@nv\endcsname{\def\PYG@tc##1{\textcolor[rgb]{0.73,0.38,0.84}{##1}}}
\expandafter\def\csname PYG@tok@s1\endcsname{\def\PYG@tc##1{\textcolor[rgb]{0.25,0.44,0.63}{##1}}}
\expandafter\def\csname PYG@tok@gp\endcsname{\let\PYG@bf=\textbf\def\PYG@tc##1{\textcolor[rgb]{0.78,0.36,0.04}{##1}}}
\expandafter\def\csname PYG@tok@sh\endcsname{\def\PYG@tc##1{\textcolor[rgb]{0.25,0.44,0.63}{##1}}}
\expandafter\def\csname PYG@tok@ow\endcsname{\let\PYG@bf=\textbf\def\PYG@tc##1{\textcolor[rgb]{0.00,0.44,0.13}{##1}}}
\expandafter\def\csname PYG@tok@sx\endcsname{\def\PYG@tc##1{\textcolor[rgb]{0.78,0.36,0.04}{##1}}}
\expandafter\def\csname PYG@tok@bp\endcsname{\def\PYG@tc##1{\textcolor[rgb]{0.00,0.44,0.13}{##1}}}
\expandafter\def\csname PYG@tok@c1\endcsname{\let\PYG@it=\textit\def\PYG@tc##1{\textcolor[rgb]{0.25,0.50,0.56}{##1}}}
\expandafter\def\csname PYG@tok@kc\endcsname{\let\PYG@bf=\textbf\def\PYG@tc##1{\textcolor[rgb]{0.00,0.44,0.13}{##1}}}
\expandafter\def\csname PYG@tok@c\endcsname{\let\PYG@it=\textit\def\PYG@tc##1{\textcolor[rgb]{0.25,0.50,0.56}{##1}}}
\expandafter\def\csname PYG@tok@mf\endcsname{\def\PYG@tc##1{\textcolor[rgb]{0.13,0.50,0.31}{##1}}}
\expandafter\def\csname PYG@tok@err\endcsname{\def\PYG@bc##1{\setlength{\fboxsep}{0pt}\fcolorbox[rgb]{1.00,0.00,0.00}{1,1,1}{\strut ##1}}}
\expandafter\def\csname PYG@tok@kd\endcsname{\let\PYG@bf=\textbf\def\PYG@tc##1{\textcolor[rgb]{0.00,0.44,0.13}{##1}}}
\expandafter\def\csname PYG@tok@ss\endcsname{\def\PYG@tc##1{\textcolor[rgb]{0.32,0.47,0.09}{##1}}}
\expandafter\def\csname PYG@tok@sr\endcsname{\def\PYG@tc##1{\textcolor[rgb]{0.14,0.33,0.53}{##1}}}
\expandafter\def\csname PYG@tok@mo\endcsname{\def\PYG@tc##1{\textcolor[rgb]{0.13,0.50,0.31}{##1}}}
\expandafter\def\csname PYG@tok@mi\endcsname{\def\PYG@tc##1{\textcolor[rgb]{0.13,0.50,0.31}{##1}}}
\expandafter\def\csname PYG@tok@kn\endcsname{\let\PYG@bf=\textbf\def\PYG@tc##1{\textcolor[rgb]{0.00,0.44,0.13}{##1}}}
\expandafter\def\csname PYG@tok@o\endcsname{\def\PYG@tc##1{\textcolor[rgb]{0.40,0.40,0.40}{##1}}}
\expandafter\def\csname PYG@tok@kr\endcsname{\let\PYG@bf=\textbf\def\PYG@tc##1{\textcolor[rgb]{0.00,0.44,0.13}{##1}}}
\expandafter\def\csname PYG@tok@s\endcsname{\def\PYG@tc##1{\textcolor[rgb]{0.25,0.44,0.63}{##1}}}
\expandafter\def\csname PYG@tok@kp\endcsname{\def\PYG@tc##1{\textcolor[rgb]{0.00,0.44,0.13}{##1}}}
\expandafter\def\csname PYG@tok@w\endcsname{\def\PYG@tc##1{\textcolor[rgb]{0.73,0.73,0.73}{##1}}}
\expandafter\def\csname PYG@tok@kt\endcsname{\def\PYG@tc##1{\textcolor[rgb]{0.56,0.13,0.00}{##1}}}
\expandafter\def\csname PYG@tok@sc\endcsname{\def\PYG@tc##1{\textcolor[rgb]{0.25,0.44,0.63}{##1}}}
\expandafter\def\csname PYG@tok@sb\endcsname{\def\PYG@tc##1{\textcolor[rgb]{0.25,0.44,0.63}{##1}}}
\expandafter\def\csname PYG@tok@k\endcsname{\let\PYG@bf=\textbf\def\PYG@tc##1{\textcolor[rgb]{0.00,0.44,0.13}{##1}}}
\expandafter\def\csname PYG@tok@se\endcsname{\let\PYG@bf=\textbf\def\PYG@tc##1{\textcolor[rgb]{0.25,0.44,0.63}{##1}}}
\expandafter\def\csname PYG@tok@sd\endcsname{\let\PYG@it=\textit\def\PYG@tc##1{\textcolor[rgb]{0.25,0.44,0.63}{##1}}}

\def\PYGZbs{\char`\\}
\def\PYGZus{\char`\_}
\def\PYGZob{\char`\{}
\def\PYGZcb{\char`\}}
\def\PYGZca{\char`\^}
\def\PYGZam{\char`\&}
\def\PYGZlt{\char`\<}
\def\PYGZgt{\char`\>}
\def\PYGZsh{\char`\#}
\def\PYGZpc{\char`\%}
\def\PYGZdl{\char`\$}
\def\PYGZhy{\char`\-}
\def\PYGZsq{\char`\'}
\def\PYGZdq{\char`\"}
\def\PYGZti{\char`\~}
% for compatibility with earlier versions
\def\PYGZat{@}
\def\PYGZlb{[}
\def\PYGZrb{]}
\makeatother

\begin{document}

\maketitle
\tableofcontents
\phantomsection\label{index::doc}


Some notes for quantum


\chapter{Introduction}
\label{index:introduction}\label{index:quantum-notes}
Some notes continued from the full theoretical physics notes are \href{http://cosmologytaskforce.github.io/PhysicsResearchSurvivalManual/}{here}.


\chapter{Table of Contents}
\label{index:table-of-contents}

\section{Vocabulary}
\label{vocabulary::doc}\label{vocabulary:vocabulary}
Vocabulary of physics, the fountain of research ideas.
\begin{enumerate}
\setcounter{enumi}{-1}
\item {} 
Fine Structure Constant

\end{enumerate}

$\alpha = \frac{k_\mathrm{e} e^2}{\hbar c} = \frac{1}{(4 \pi \varepsilon_0)} \frac{e^2}{\hbar c} = \frac{e^2 c \mu_0}{2 h}$

In electrostatic cgs units, $\alpha = \frac{e^2}{\hbar c}$.

In natural units, $\alpha = \frac{e^2}{4 \pi}$ .
\begin{enumerate}
\item {} 
Hydrogen Atom

\end{enumerate}

Potential $V(r) = -\frac{Z e^2}{4\pi \epsilon_0 r}$.

Energy levels: $E_{n} = -\left(\frac{Z^2 \mu e^4}{32 \pi^2\epsilon_0^2\hbar^2}\right)\frac{1}{n^2} = -\left(\frac{Z^2\hbar^2}{2\mu a_{\mu}^2}\right)\frac{1}{n^2} = \frac{\mu c^2Z^2\alpha^2}{2n^2}.$

Ground state of hydrogen atom $\psi_{100}(r)=\frac{1}{\sqrt{\pi}}\frac{1}{a^{3/2}} e^{-Z r/a}$.


\section{Approximation Methods}
\label{approx::doc}\label{approx:approximation-methods}

\subsection{Variational Method}
\label{approx:variational-method}

\subsubsection{Trial functions}
\label{approx:trial-functions}
Some of the calculable trial functions:
\begin{enumerate}
\item {} 
$\psi(x) = \cos\alpha x$, for $|\alpha x|<\pi/2$, otherwise 0.

\item {} 
$\psi(x) = \alpha^2 - x^2$, for $|x|<\alpha$, otherwise 0.

\item {} 
$\psi(x) = C \exp(-\alpha x^2/2)$.

\item {} 
$\psi(x) = C(\alpha - |x|)$, for $|x|<\alpha$, otherwise 0.

\item {} 
$\psi(x) = C\sin\alpha x$, for $|\alpha x|<\pi$, otherwise 0.

\end{enumerate}


\subsubsection{Procedure}
\label{approx:procedure}\begin{enumerate}
\item {} 
Pick a trial function.

\begin{notice}{note}{Note:}
How to pick a trial function? For ground state energy, we should pick a function that has the same property as the real ground state. This requires some understanding of the problem we are dealing with.

Things to consider:
\begin{enumerate}
\item {} 
The new problem is just a modification of a known solved problem. Then we can easily find out what really is different and interprete the new problem in terms of the old one.

\item {} 
If the Hamiltonian have definite parity, the ground state wave function should pick up some parity which is usually even to make it the lowest energy.

\item {} 
Continious function? A $C^\infty$ Hamiltonian can only have continious functions as solutions for a finite system.

\item {} 
Nodes deteremines the kinetic energy so check the nodes for ground state wave function.

\item {} 
Check the behivior of the wave function at different limits. In most cases, the Shrödinger equation can be reduced to something solvable at some limits.

\item {} 
\textbf{One more thing, the trial function should make the problem calculable.}

\end{enumerate}
\end{notice}

\end{enumerate}


\subsubsection{Why Not General Viriational Method}
\label{approx:why-not-general-viriational-method}
Why don't we just use a most general variational method to find out the ground state? Because we will eventually come back to the time-independent Shrödinger equation.

Suppose we have a functional form
\begin{gather}
\begin{split}E(\psi^*, \psi, \lambda) = \int dx \psi^* H \psi - \lambda \left( \int dx \psi^* \psi - 1 \right)\end{split}\notag\\\begin{split}\end{split}\notag
\end{gather}
The reason we have this Lagrange multiplier method is that the wave function should be normalized and this multiplier provides the degree of freedom. We would only get a wrong result if we don't include this DoF.

Variation of $\psi^*$,
\begin{gather}
\begin{split}\delta E = \int dx \delta \psi^* H \psi - \int dx \delta \psi^* \psi = 0\end{split}\notag\\\begin{split}\end{split}\notag
\end{gather}
Now what?
\begin{gather}
\begin{split}H \psi - \lambda \psi = 0\end{split}\notag\\\begin{split}\end{split}\notag
\end{gather}
Not helpful.


\subsection{Variational Method and Virial Theorem}
\label{approx:variational-method-and-virial-theorem}
For a potential $V(x)=b x^n$, we can prove that virial theorem is valid for ground state if we use Gaussian trial function $e^{- \alpha x^2/2}$.

A MMA proof is here.

Virial theorem is pretty interesting. It shares the same math with equipartition theorem.


\subsection{WKB}
\label{approx:wkb}
This is a semi-classical method. It is semi classical because we will use the classical momentum
\begin{gather}
\begin{split}\hbar k(x) = \sqrt{2m (E - V(x))}\end{split}\notag\\\begin{split}\end{split}\notag
\end{gather}
The following points are important for this method.
\begin{enumerate}
\setcounter{enumi}{-1}
\item {} 
WKB start from a classical estimation of wave number at a certain energy $E$ which is later quantified by the Bohr-Sommerfeld quantization rule.

\item {} 
Conservation law:
\begin{gather}
\begin{split}\frac{\partial}{\partial t}\rho + \nabla \cdot \vec j = 0\end{split}\notag\\\begin{split}\end{split}\notag
\end{gather}
where $\rho = \psi^* \psi$, $\vec j = -\frac{\hbar}{2 m i} \left( \psi \nabla \psi^* - \psi^* \nabla \psi \right)$. This can be derived from Shrödinger equation easily.

\item {} 
Phase:
Wave function is generally $A(x)\exp(\phi(x))$. However, $\phi(x)$ should be the area of the phase function starting from some initial point. For example in WKB, $k(x) = \phi'(x)$ and $\phi(x) = \int \phi'(x')d x' = \int k(x') d x'$.

Using this general wave function and conservation law we find out that $A(x) ~ \frac{1}{\sqrt{k(x)}}$. Then we can apply the two boundary conditions. However we will find two different wave functions given by two boundary conditions. Now we should connect them because $\psi(a) = \psi(b)$ exactly. By comparing the two wave functions we can find something like Bohr-Sommerfeld quantization rule.

\item {} 
Correction at bouldary:
However, this method requires that the potential varies slowly or equivalently the wave number varies slowly. Basicly we are just using the following approximation:
\begin{gather}
\begin{split}A'(x) = 0, k'(x) = 0\end{split}\notag\\\begin{split}\end{split}\notag
\end{gather}
For example when taking the derivative of wave function,
\begin{gather}
\begin{split}\psi'(x) = A'(x) e^{i\int \cdots} + A(x) k(x) e^{i\int \cdots} \approx A(x) k(x) e^{i\int \cdots}\end{split}\notag\\\begin{split}\end{split}\notag
\end{gather}
where we drop the term with $A'(x)$. That is to say
\begin{gather}
\begin{split}|A'|\ll |A k| \Rightarrow |k'| \ll k^2\end{split}\notag\\\begin{split}\end{split}\notag
\end{gather}
But at boundary where $E = V$, this is obviously not valid because $k=0$. So we need to fix this problem.

The solution is to use first order of the potential in a Taylor expansion. Then solve the problem exactly. Finally we connect regions that is far out from the boundary, need the boundary and between the boundary.

\end{enumerate}

If we can have a good boundary condition, then the energy spectrum given by WKB can be very good. Even we don't have a good boundary condtion, the excited states given by this method are always close to the exact ones.


\bigskip\hrule{}\bigskip



\bigskip\hrule{}\bigskip

\begin{figure}[htbp]
\centering
\href{http://creativecommons.org/licenses/by-nc-sa/3.0/us/}{\includegraphics{cc_byncsa.png}}\end{figure}


\bigskip\hrule{}\bigskip


This open source project is hosted on GitHub: \href{https://github.com/emptymalei/quantum}{quantum} .

\href{https://raw.github.com/emptymalei/quantum/master/Notes/\_build/latex/quantum.pdf}{Latest PDF here}.



\renewcommand{\indexname}{Index}
\printindex
\end{document}
